% =====================================================================
\section{Gray-map compatibility and finite-field representation}
\label{sec:gray}
% =====================================================================

This section is the bridge from the ring $R$ to $\F_q$. The map is the
coefficient expansion; the content is the identification of the form it
transports.

\subsection{The coefficient expansion}

\begin{definition}[coefficient expansion]\label{def:gray}
For $a=\sum_{a'=0}^{t-1}c_{a'}u^{a'}\in A_t$ put
$\varphi_A(a)=(c_0,\dots,c_{t-1})\in\F_q^{\,t}$. For
$X=(X_{kl})\in R=M_2(A_t)$ put
$\varphi_R(X)=(\varphi_A(X_{00}),\varphi_A(X_{01}),\varphi_A(X_{10}),
\varphi_A(X_{11}))\in\F_q^{\,4t}$, and for $x=(x_0,\dots,x_{n-1})\in R^n$,
\[
  \varphi_n(x)=\bigl(\varphi_R(x_0),\dots,\varphi_R(x_{n-1})\bigr)
  \in\F_q^{\,\Nlen},\qquad \Nlen=4tn .
\]
The flat index is $\varphi_n(x)_{(e\cdot4+k)t+a}=
\text{the $u^a$-coefficient of the $k$-th entry of }x_e$.
\end{definition}

\begin{lemma}\label{lem:philin}
$\varphi_n:R^n\to\F_q^{\Nlen}$ is an $\F_q$-linear bijection; it is additive
and satisfies $\varphi_nT=\widehat T\varphi_n$, where $T$ is the cyclic shift
on $R^n$ and $\widehat T$ is the rotation by $4t$ on $\F_q^{\Nlen}$. It is
\textbf{not} a ring homomorphism: $R$ is not commutative, and
$\varphi_R(XY)\ne\varphi_R(YX)$ in general.
\end{lemma}

\begin{proof}
$\varphi_A$ is a bijection $A_t\to\F_q^t$ by the uniqueness of the normal form;
$\varphi_R$ and $\varphi_n$ are products of it, hence bijections, and they are
$\F_q$-linear because taking the $u^a$-coefficient is. For the shift:
$T$ sends $x_e$ to $x_{e-1}$, and $\widehat T$ rotates the flat index by $4t$,
which is exactly one $R$-block; the two agree. Non-multiplicativity is the
explicit computation $\varphi_R(XY)\ne\varphi_R(YX)$ on matrix units.
\end{proof}

\subsection{The exact induced form}\label{subsec:induced}

Because $\varphi_n$ is an $\F_q$-linear bijection, there is exactly one form on
$\F_q^{\Nlen}$ for which $\varphi_n$ is a form-isometry, namely the transport of
$\bb_s$. The content of the following theorem is to identify it.

\begin{theorem}[Gram matrix of the transported form]\label{thm:gram}
Transport $\bb_s$ along $\varphi_n$:
$\langle u,v\rangle^{\widehat{}}_s:=\bb_s(\varphi_n^{-1}u,\varphi_n^{-1}v)$.
Then, writing $u^{[s]}=(\sigma_s(u_i))_i$,
\[
  \langle u,v\rangle^{\widehat{}}_s=u^{[s]}\,J\,v^{\mathsf T},
  \qquad
  J=I_{4n}\otimes H_t,\qquad
  (H_t)_{ab}=\begin{cases}1,&a+b=t-1,\\0,&\text{else.}\end{cases}
\]
The matrix $J$ is non-singular, symmetric and Hankel, $J^2=I$,
$J=I$ iff $t=1$, and $J$ is alternating iff $t$ is even.
\end{theorem}

\begin{proof}
First, $\tr(\theta_s(X)Y)=\sum_{k,l}\sigma_s(X_{kl})Y_{kl}$: indeed
$\theta_s(X)=\sigma_s(X^{\mathsf T})$ has $(k,l)$-entry
$\sigma_s(X_{lk})$, so
$\tr(\theta_s(X)Y)=\sum_{k,l}\sigma_s(X_{lk})Y_{lk}$, and
$(k,l)\mapsto(l,k)$ is a bijection of the index set.

Therefore
\[
  \bb_s(x,y)=\sum_e\sum_{k,l}
  \lambda\bigl(\sigma_s((x_e)_{kl})(y_e)_{kl}\bigr)
  =\sum_{e,k}\sum_{a=0}^{t-1}
  \sigma_s\bigl((x_e)_{k,a}\bigr)\,(y_e)_{k,t-1-a},
\]
because $\lambda(u^{a+b})=1$ iff $a+b=t-1$ and $0$ otherwise, where
$(x_e)_{k,a}$ denotes the $u^a$-coefficient of the $k$-th entry of $x_e$.

Reading the right-hand side in the flat coordinates
$u=\varphi_n(x)$, $v=\varphi_n(y)$, the summand pairs the coordinate
$(e,k,a)$ of $u$ (with $\sigma_s$ applied) with the coordinate
$(e,k,t-1-a)$ of $v$. In the ordering of Definition~\ref{def:gray} the
$a$-index runs inside a fixed $(e,k)$ block, so the Gram matrix is the identity
on the $(e,k)$-index tensored with $H_t$ on the $a$-index, i.e.
$J=I_{4n}\otimes H_t$.

The stated properties of $J$: $H_t$ is the reversal (Hankel) matrix, symmetric
with $H_t^2=I_t$; hence $J$ is symmetric with $J^2=I_{4n}\otimes H_t^2=I$,
non-singular. $J=I$ iff $H_t=I_t$, i.e. $t=1$. Finally $H_t$ is alternating
(zero diagonal, since $2a=t-1$ has no solution when $t$ is even) exactly when
$t$ is even, and $J$ inherits this.
\end{proof}

\begin{corollary}[the transported form is not Euclidean]\label{cor:noteuclid}
The transported form is Euclidean ($\langle u,v\rangle^{\widehat{}}_s=u\cdot v$)
if and only if $t=1$ and $r\mid s$; it is the plain
$\widehat\sigma_s$-Hermitian form if and only if $t=1$.
\end{corollary}

\begin{proof}
The Gram matrix of the Euclidean form is $I$; by
Theorem~\ref{thm:gram} the Gram matrix is $I_{4n}\otimes H_t$ with the
$\sigma_s$ twist, which equals $I$ iff $H_t=I_t$ (i.e. $t=1$) and
$\sigma_s=\id$ (i.e. $r\mid s$). Dropping the $\sigma_s$ condition gives the
plain Hermitian statement.
\end{proof}

\subsection{Duality and hull preservation}

\begin{theorem}[duality preservation]\label{thm:grayduality}
For every code $C\le\Rn$ and every $s$,
\[
  \varphi_n(C^{\perp_s})=\varphi_n(C)^{\hhatpers},
\]
where $D^{\hhatpers}=\{v:\langle u,v\rangle^{\widehat{}}_s=0\ \forall u\in D\}$.
\end{theorem}

\begin{proof}
$\varphi_n$ is bijective and $\langle\cdot,\cdot\rangle^{\widehat{}}_s$ is
defined as its transport, so for $y\in R^n$,
\[
  y\in C^{\perp_s}
  \iff \bb_s(x,y)=0\ \forall x\in C
  \iff \langle\varphi_nx,\varphi_ny\rangle^{\widehat{}}_s=0\ \forall x\in C
  \iff \varphi_n(y)\in\varphi_n(C)^{\hhatpers}. \qedhere
\]
\end{proof}

\begin{theorem}[hull preservation, with no principality hypothesis]
\label{thm:grayhull}
For every $C\le\Rn$ and every $s$,
\[
  \varphi_n\bigl(\Hull_s(C)\bigr)
  =\varphi_n(C)\cap\varphi_n(C)^{\hhatpers}.
\]
In particular the identity holds for the non-principal hulls of
Theorem~\ref{thm:nonprincipal}.
\end{theorem}

\begin{proof}
$\varphi_n$ is injective, hence preserves intersections:
$\varphi_n(C\cap C^{\perp_s})=\varphi_n(C)\cap\varphi_n(C^{\perp_s})$; now
apply Theorem~\ref{thm:grayduality}. No principality hypothesis is used at any
point: Definition~\ref{def:hull} is an intersection of $\F_q$-subspaces and
Theorems~\ref{thm:hullstruct} and~\ref{thm:grayduality} hold for arbitrary right
$\Rn$-submodules.
\end{proof}

\subsection{The image code}

\begin{proposition}[quasi-cyclicity]\label{prop:qc}
Let $D=\varphi_n(C)\subseteq\F_q^{\Nlen}$. Then $D$ is $\F_q$-linear of length
$\Nlen$ and dimension $\dim_{\F_q}C$, and it is \emph{quasi-cyclic of index
$4t$ and co-index $n$}: it is invariant under the rotation by $4t$
positions. It is not cyclic in general.
\end{proposition}

\begin{proof}
Linearity and dimension are Lemma~\ref{lem:philin}. Invariance under the
rotation by $4t$ is the shift identity $\varphi_nT=\widehat T\varphi_n$
together with $T$-invariance of the right ideal $C$. Non-cyclicity: the
rotation by one position does not in general preserve $D$; this is a
measurement on the enumerated families and we record it rather than claim it.
\end{proof}

\begin{remark}[what the map does to weights]\label{rem:weights}
Let $w_{\rm H}$ denote the Hamming weight on $R^n$ (counting nonzero
$R$-positions), $w_{\rm ent}$ the weight counting nonzero entries of the
$2\times2$ matrices, and $w_{\rm coef}$ the number of nonzero $\F_q$-coefficients.
Then
\[
  w_{\rm H}\le w_{\rm ent}\le w_{\rm coef}\le t\,w_{\rm ent},
  \qquad
  w_{\rm H}\le\wt_{\rm H}\varphi_n\le 4t\,w_{\rm H},
\]
and both bounds are sharp. In particular $\varphi_n$ \emph{is} an isometry from
the coefficient metric $d_{\rm coef}$ to the Hamming metric on
$\F_q^{\Nlen}$. The naive claim that $\varphi_n$ preserves the ring-Hamming
weight is false: a nonzero matrix with a single nonzero entry has
$w_{\rm H}=1$ but $\wt_{\rm H}\varphi_n\ge1$ with equality only in special
cases, and in general the two differ.
\end{remark}

\begin{remark}[no homogeneous-weight isometry]\label{rem:noisometry}
The homogeneous weight on $R=M_2(\F_q)$ takes the values $0$,
$1+\frac1{q^2-1}$ (rank one) and $1-\frac1{(q-1)^2(q+1)}$ (units). Since these
values are not integers, \emph{no} map into any Hamming space can be an isometry
for it. The chain of rings for which such isometries do exist
\cite{GreferathSchmidt1999GrayIsometry} consists of finite \emph{chain} rings,
and $R$ is not one. This places a hard limit on what any Gray map can achieve
here, and it is the reason Section~\ref{subsec:phys} recomputes distances in
physical coordinates rather than transporting them.
\end{remark}

\subsection{Every Galois hull is self-orthogonal}

This is the hinge of the paper. It holds in \emph{all three regimes} of
Definition~\ref{def:regimes}, not only the reflexive one.

\begin{theorem}[hull self-orthogonality]\label{thm:hullso}
For every $C\le\Rn$ and every $s$,
\[
  \varphi_n\bigl(\Hull_s(C)\bigr)
  \subseteq
  \varphi_n\bigl(\Hull_s(C)\bigr)^{\hhatpers}.
\]
No reflexivity hypothesis ($r\mid 2s$) and no principality hypothesis is
required.
\end{theorem}

\begin{proof}
Let $H=\Hull_s(C)=C\cap C^{\perp_s}$ and let $x,y\in H$. Then $x\in C$ and
$y\in C^{\perp_s}$. By the definition of $C^{\perp_s}$ we have
$\bb_s(c,y)=0$ for every $c\in C$; taking $c=x$ gives $\bb_s(x,y)=0$.
Transporting with $\varphi_n$ (Theorem~\ref{thm:gram}) gives
$\langle\varphi_nx,\varphi_ny\rangle^{\widehat{}}_s=0$.

Note what the argument uses and what it does not: it uses $x\in C$ and
$y\in C^{\perp_s}$, and it never interchanges the two slots. It therefore does
not need the orthogonality relation to be symmetric, which is exactly the
property that holds only in the reflexive regime
(Proposition~\ref{prop:regimes}). Nor does it use principality: $H$ is an
intersection of $\F_q$-subspaces, and Theorems~\ref{thm:hullstruct}
and~\ref{thm:grayduality} hold for arbitrary right $\Rn$-submodules.

Finally, note the direction of the statement: a nonzero hull dimension never
implies that $C$ itself is self-orthogonal.
\end{proof}

\subsection{Physical coordinates and the characteristic-two obstruction}
\label{subsec:phys}

For the Euclidean (CSS) route one wants the transported form to be the Euclidean
form. Corollary~\ref{cor:noteuclid} says this happens only for $t=1$ with
$r\mid s$; otherwise one looks for a linear change of coordinates.

\begin{definition}\label{def:congruence}
Assume $\sigma_s=\id$ (i.e. $r\mid s$), so that the transported form is the
bilinear form with matrix $J$. A \emph{congruence} is a matrix $\mathcal T$
with $\mathcal T\mathcal T^{\mathsf T}=J$; the associated coordinate map is
$\Psi(u)=u\,\mathcal T$.
\end{definition}

\begin{proposition}[the congruence and its cost]\label{prop:congruence}
Let $T$ be a $t\times t$ matrix with $T^{\mathsf T}H_tT=I_t$ and let
$\mathcal T=I_{4n}\otimes T$. Then $\mathcal T\mathcal T^{\mathsf T}=J$ (using
$J^2=I$ and $H_t^{\mathsf T}=H_t$), $\mathcal T$ is invertible, and $\Psi$
transports the transported form to the Euclidean form:
\[
  \Psi(u)\cdot\Psi(v)=u\,\mathcal T\mathcal T^{\mathsf T}v^{\mathsf T}
  =u\,J\,v^{\mathsf T}=\langle u,v\rangle^{\widehat{}}_s .
\]
Moreover $\mathcal T$ commutes with the block shift (so the quasi-cyclic
structure of index $4t$ is preserved), and $\mathcal T$ is monomial iff $t=1$.
\end{proposition}

\begin{proof}
Immediate from $J=I_{4n}\otimes H_t$, $H_t^{\mathsf T}=H_t$, $J^2=I$ and the
definitions; $\mathcal T$ is block-diagonal in the flat ordering, hence commutes
with the block shift. Monomiality: for $t=1$, $H_1=(1)$ and $T$ may be taken to
be $(1)$; for $t\ge2$, $T^{\mathsf T}H_tT=I_t$ forces $T$ to have at least two
nonzero entries in some row, so $\mathcal T$ is not monomial.
\end{proof}

\begin{theorem}[characteristic-two obstruction]\label{thm:obstruction}
Let $p=2$ and let $t$ be even. Then no matrix $T$ with
$T^{\mathsf T}H_tT=I_t$ exists. Consequently, for every such $(q,t)$ and
\emph{every} code, the Euclidean/CSS route is not available: no linear change of
coordinates converts the transported form into the Euclidean form.
\end{theorem}

\begin{proof}
If $t$ is even then $H_t$ is alternating: its diagonal entries are
$(H_t)_{aa}=1$ iff $2a=t-1$, which has no integer solution when $t$ is even;
and $H_t$ is symmetric. A symmetric matrix congruent to the identity over a
field of characteristic $2$ cannot be alternating, because an alternating
matrix $M$ satisfies $x^{\mathsf T}Mx=0$ for all $x$ whereas
$x^{\mathsf T}Ix=|x|^2$ is not identically zero. Hence no such $T$ exists.
\end{proof}

\begin{remark}[the hypotheses are exact]\label{rem:obstructionexact}
Theorem~\ref{thm:obstruction} is stated for $p=2$ \emph{and} $t$ even, and it
must not be extended. Congruences do exist for $t=1$, $t=3$ and $t=5$ at
$q=4,8$, and one exists at $(q,t)=(9,2)$. What the obstruction forbids is
exactly the pair ``characteristic two, $t$ even''; it is a statement about the
form $J$, not about any code. Where a congruence exists but is non-monomial
(Proposition~\ref{prop:congruence}), it must not be treated as an isometry:
measured instances show that the minimum distance of a code \emph{does} change
under $\Psi$ (3 of 5 codes at $(9,2,2)$), so every distance must be recomputed
in physical coordinates.
\end{remark}

\begin{theorem}[number of physical coordinates]\label{thm:nphys}
A quantum code produced from a code $C\le\Rn$ by any of the routes of
Section~\ref{sec:quantum} has length $\Nlen=4tn$ over the alphabet $\F_q$.
\end{theorem}

\begin{proof}
One qudit of dimension $q$ carries one $\F_q$-coordinate, so the number of
physical qudits equals the $\F_q$-length of the classical code whose symplectic
group lives over $\F_q^{2\Nlen}$: by Lemma~\ref{lem:basic}(c),
$\dim_{\F_q}R^n=\Nlen=4tn$. There is no doubling to $2\Nlen$: $\F_q^{2\Nlen}$
is the symplectic \emph{bookkeeping} space in which the stabilizer group sits,
not a set of qudits. Nor is there a substitution of $\F_q$ by $\F_p$: the
alphabet of the quantum code is $\F_q$ throughout.
\end{proof}

\subsection{Worked examples IV: the Gray map}

\begin{example}[the coefficient expansion and the Gram matrix, at $t=2$]
\label{ex:gray}
Let $q=4$, $t=2$, $n=1$, so $\Nlen=4\cdot2\cdot1=8$. Take
\[
  X=\begin{pmatrix}1+\alpha u&\beta u\\ 0&1\end{pmatrix}\in R=M_2(A_2),
  \qquad \alpha,\beta\in\F_4 .
\]
Then $\varphi_A(1+\alpha u)=(1,\alpha)$, $\varphi_A(\beta u)=(0,\beta)$,
$\varphi_A(0)=(0,0)$, $\varphi_A(1)=(1,0)$, so
\[
  \varphi_R(X)=\bigl(1,\alpha,\;0,\beta,\;0,0,\;1,0\bigr)\in\F_4^8 ,
\]
in the flat ordering of Definition~\ref{def:gray}: the four $2\times2$ entries
in reading order, each expanded into its $t=2$ coefficients.

For $t=2$ the Hankel matrix is
\[
  H_2=\begin{pmatrix}0&1\\1&0\end{pmatrix},
  \qquad\text{so}\qquad
  J=I_{4n}\otimes H_2=I_4\otimes H_2 ,
\]
an $8\times8$ matrix which is symmetric, satisfies $J^2=I_8$, is \emph{not}
the identity (in accordance with Corollary~\ref{cor:noteuclid}, since
$t\ne1$), and \emph{is} alternating, since $(H_2)_{aa}=1$ would require
$2a=1$ and $p=2$. The transported form is therefore
\[
  \langle u,v\rangle^{\widehat{}}_s=u^{[s]}Jv^{\mathsf T}
  =\sum_{i=0}^{3}\bigl(\sigma_s(u_{2i})v_{2i+1}
    +\sigma_s(u_{2i+1})v_{2i}\bigr),
\]
which pairs the two $u$-coefficients inside each of the four matrix entries
and nothing else. Note that it is not the dot product $u\cdot v$: the
coordinate $u_0$ is paired with $v_1$, not with $v_0$.
\end{example}

\begin{example}[the characteristic-two obstruction, at $(q,t)=(4,2)$]
\label{ex:obstruction}
Let $q=4$ (so $p=2$) and $t=2$. Then $H_2=\begin{pmatrix}0&1\\1&0\end{pmatrix}$
has zero diagonal, and for every $x=(x_0,x_1)^{\mathsf T}\in\F_4^2$,
\[
  x^{\mathsf T}H_2x=x_0x_1+x_1x_0=2x_0x_1=0
\]
in characteristic $2$: $H_2$ is alternating. For any $2\times2$ matrix $T$ over
$\F_4$ and any $x$,
\[
  x^{\mathsf T}(T^{\mathsf T}H_2T)x=(Tx)^{\mathsf T}H_2(Tx)=0,
\]
so $T^{\mathsf T}H_2T$ is alternating for \emph{every} $T$. But $I_2$ is not
alternating, since $e_1^{\mathsf T}I_2e_1=1\ne0$. Hence no $T$ with
$T^{\mathsf T}H_2T=I_2$ exists; the search over all $4^4=256$ matrices
$T\in M_2(\F_4)$ confirms this exhaustively, and Theorem~\ref{thm:obstruction}
proves it for every even $t$ in characteristic $2$.

The hypotheses are exact, and the neighbouring cases are different: for
$t=1,3,5$ at $q=4,8$ a congruence is \emph{constructed} (symmetric elimination
with Frobenius scaling, backtracking over pivots) and verified by direct
multiplication, and one exists at $(q,t)=(9,2)$ as well --- there $p=3$ and the
argument above breaks because $2x_0x_1$ need not vanish. Where a congruence
exists but is non-monomial (Proposition~\ref{prop:congruence}), it must not
be treated as an isometry: at $(9,2,2)$ the minimum distance changed under
$\Psi$ for $3$ of the $5$ codes tested, so every distance in this paper is
recomputed in physical coordinates.
\end{example}
